% Propietario conservado: /Users/ruben/Documents/New project/output/REVISION_ARGUMENTAL_APERTURAS_CIERRES_20260915/VI/delivery/MANUSCRIPT_VI_ES_EN_COMPLETE/source_es/sections/03_conjugacion_mobius.tex
\section{Conjugaison, orientation et retour des sections matérielles}
\label{vi:vi:sec:conjugacion}

L’identité persistante de la section~\ref{vi:vi:sec:particula} admet
des opérations d’inversion à différents niveaux de sa structure.
L’inversion cellulaire agit dans l’atlas ; la conjugaison dodécaphasique
agit dans le registre d’orientation ; le relèvement d’une boucle agit
dans une fibre ; une rotation spinorielle agit dans un espace de
représentation. Spécifier ces domaines permet de composer les
opérations et d’expliquer quelles informations sont conservées à chaque retour.

\subsection{Involution dodécaphasique et parité des observables}
\label{vi:vi:subsec:conjugacion-registro}

Soit $\mathcal T_{12}=\{-1,0,1\}^{12}$ l’espace ambiant des registres
tritiques de douze positions. Le domaine admissible du TPK est le
sous-ensemble muni des compatibilités de route déjà construites. L’opération
sur l’espace ambiant est
\begin{equation}
 C_M:\mathcal T_{12}\longrightarrow\mathcal T_{12},\qquad
 (C_M\tau)_j=-\tau_{11-j},\quad 0\leq j<12.
 \label{vi:vi:eq:cm}
\end{equation}
La restriction aux registres admissibles utilise la stabilité de ce
domaine sous la conjugaison. L’ordre inversé et le changement de signe
sont des composantes simultanées de l’opération.

\begin{proposition}[Parité linéaire et quadratique du registre]
\label{vi:vi:prop:paridad-registro}
On a $C_M^2=I$. Si $a_j=a_{11-j}$ et
$L_a(\tau)=\sum_ja_j\tau_j$, alors
$L_a(C_M\tau)=-L_a(\tau)$. Si une forme quadratique
$Q_b(\tau)=\sum_{j,k}b_{jk}\tau_j\tau_k$ satisfait
$b_{jk}=b_{11-j,11-k}$, alors $Q_b(C_M\tau)=Q_b(\tau)$.
\end{proposition}
\begin{proof}
À chaque position,
$[C_M(C_M\tau)]_j=-(-\tau_{11-(11-j)})=\tau_j$.
Pour la fonctionnelle linéaire, substituer $k=11-j$ dans la somme donne
$L_a(C_M\tau)=-\sum_ka_{11-k}\tau_k=-L_a(\tau)$.
Dans l’expression quadratique, les deux changements de signe s’annulent et
la substitution $(r,s)=(11-j,11-k)$ récupère $Q_b$ par la
symétrie de ses coefficients.
\end{proof}

L’orientation linéaire et les corrélations de mémoire ont donc
des lois différentes. Même une composition de ces observables doit
conserver leurs parités : si $L_a$ est impair, son caractère exponentiel
satisfait
\begin{equation}
 \exp(L_a(C_M\tau))=\exp(L_a(\tau))^{-1}.
 \label{vi:vi:eq:caracter-inverso}
\end{equation}
L’égalité exprime une inversion multiplicative. L’invariance d’un
opérateur matériel complet s’étudie sur cet opérateur, avec ses
canaux et sa représentation, et se démontre par la covariance
correspondante.

L’inversion cellulaire du théorème~\ref{vi:vi:thm:selector-central} est
$\rho_c:\mathcal C_{81}\to\mathcal C_{81}$, tandis que $C_M$ a pour
domaine \eqref{vi:vi:eq:cm}. Toutes deux sont des involutions, mais pour
les comparer, on requiert l’application de registre de la route et sa loi de
transport. Cette donnée évite de remplacer la conjugaison d’une histoire
par l’inversion d’une coordonnée isolée de l’atlas.

\subsection{Relèvement aux antisections et conservation de la masse}
\label{vi:vi:subsec:antisecciones}

La conjugaison s’applique à la particule au moyen de deux applications
compatibles. L’application de base $C_m^{\mathrm b}$ transporte l’état
$x_m$ vers son état conjugué $\bar x_m$ ; le relèvement
$\widetilde C_m:\mathscr E_m(x)\to\mathscr E_m(\bar x)$
transporte ses sections matérielles. Soient $t_m$ le lecteur du
registre dodécaphasique et $\bar p_{n,m}$ les restrictions dans la
famille conjuguée. Leurs conditions sont la conservation de
l’admissibilité, l’involutivité et la compatibilité :
\begin{equation}
 \begin{aligned}
 \rho_{n,m}C_n^{\mathrm b}&=C_m^{\mathrm b}\rho_{n,m},
 &t_m(C_m^{\mathrm b}x_m)&=C_Mt_m(x_m),\\
 \bar p_{n,m}\widetilde C_n&=\widetilde C_mp_{n,m},
 &\widetilde C_{m,\bar x}\widetilde C_{m,x}&=I.
 \end{aligned}
 \label{vi:vi:eq:levantamiento-cm}
\end{equation}
Dans la première ligne, le lecteur publie le registre d’orientation
de l’histoire. Le reste du transport de l’état accompagne cette
publication. La seconde ligne donne les lois du relèvement ;
les indices $x,\bar x$ explicitent ses deux fibres de départ.
Si $s$ est persistante, son antisection est
$\bar s=(\widetilde C_ms_m)_m$.

\begin{proposition}[Persistance de l’antisection]
\label{vi:vi:prop:antiseccion-persistente}
Les conditions \eqref{vi:vi:eq:levantamiento-cm} transportent les sections
persistantes vers des sections persistantes et restituent la section lorsqu’on
conjugue deux fois. Si elles transportent en outre les symétries internes par
$g_m\mapsto\widetilde C_mg_m\widetilde C_m^{-1}$, elles induisent une
involution sur les classes matérielles de la
définition~\ref{vi:vi:def:particula}.
\end{proposition}
\begin{proof}
La compatibilité donne
$\bar p_{n,m}\bar s_n=\widetilde C_mp_{n,m}s_n
=\widetilde C_ms_m=\bar s_m$.
L’involutivité restitue $s$. Pour $s'_m=g_ms_m$, leurs
antisections sont reliées au moyen de
$\widetilde C_mg_m\widetilde C_m^{-1}$. La loi de restriction
dans \eqref{vi:vi:eq:levantamiento-cm} conserve la compatibilité de cette
famille, qui appartient aux groupes internes en vertu de la condition de
l’énoncé. La construction descend ainsi aux classes.
\end{proof}

La représentation de charge utilise la partie impaire du registre
matériel. Si $\mathfrak q$ est le caractère de charge déclaré, sa
loi est $\mathfrak q(\bar s)=-\mathfrak q(s)$. La charge élémentaire
qui intervient lorsqu’on exprime cette coordonnée dans une droite dimensionnelle est
la sortie interne déjà construite ; la représentation s’applique ensuite
à l’état, sans utiliser de charge cible pour choisir son registre.

Pour la masse, on a besoin de la covariance de son opérateur. Soient
$\widehat M_s$ et $\widehat M_{\bar s}$ les opérateurs autoadjoints
des deux réalisations, et soit $\mathcal C$ leur opérateur d’entrelacement
unitaire ou antiunitaire. Leur domaine se transporte au moyen de
$\mathcal C\operatorname{Dom}(\widehat M_s)
=\operatorname{Dom}(\widehat M_{\bar s})$. La condition opératorielle est
\begin{equation}
 \mathcal C\widehat M_s\mathcal C^{-1}
 =\widehat M_{\bar s}.
 \label{vi:vi:eq:covariancia-masa}
\end{equation}

\begin{proposition}[Conservation spectrale sous conjugaison]
\label{vi:vi:prop:masa-conjugada}
Sous \eqref{vi:vi:eq:covariancia-masa}, le transport d’une section
propre de masse $m\in\mathbb R$ conserve cette valeur propre. La mesure
spectrale se transporte au moyen de
$E_{\bar s}(B)=\mathcal C E_s(B)\mathcal C^{-1}$ pour tout
borélien réel $B$. En particulier, une réalisation scalaire conserve
la masse lors du passage à l’antisection.
\end{proposition}
\begin{proof}
Si $\widehat M_s\psi=m\psi$, la covariance donne
$\widehat M_{\bar s}\mathcal C\psi
=\mathcal C(m\psi)=m\mathcal C\psi$ ; la dernière égalité
vaut aussi dans le cas antiunitaire parce que $m$ est réel. Le calcul
fonctionnel réel des opérateurs autoadjoints transporte leurs
projecteurs spectraux et prouve la seconde affirmation.
\end{proof}

Dans un registre de bande, la même condition peut se voir sous une forme
élémentaire. Étant donnée une lecture réelle $E_0(\tau)$, soient
$E_\pm=(E_0\pm E_0\circ C_M)/2$ et soit $\xi\in\{+1,-1\}$
l’orientation de feuillet. Alors
\begin{equation}
 E_{\mathrm{band}}(\tau,\xi)=E_+(\tau)+\xi E_-(\tau),\qquad
 E_{\mathrm{band}}(C_M\tau,-\xi)=E_{\mathrm{band}}(\tau,\xi).
 \label{vi:vi:eq:banda-paridad}
\end{equation}
L’égalité s’obtient en utilisant le fait que $E_+$ est pair et $E_-$ impair.
Elle explique comment l’énergie de bande peut être conservée en transportant
simultanément registre et feuillet ; son identification à la masse d’une
réalisation concrète utilise \eqref{vi:vi:eq:covariancia-masa}.

Le registre électronique
$\tau_e=(+,+,+,-,-,-,+,+,+,-,-,-)$ satisfait $C_M\tau_e=\tau_e$.
Cette égalité concerne le mot de douze positions. La paire
$(\tau_e,\xi)$ conserve une orientation supplémentaire que la
conjugaison change. La distinction permet de maintenir le centre
électronique unique et, à la fois, les réalisations conjuguées de charge ;
un mot fixe ne devient pas une preuve de neutralité.

\subsection{Holonomie de signe et bande de persistance}
\label{vi:vi:subsec:mobius}

Soit $\ell$ une boucle de retour du lecteur matériel d’une extension
survivante. Son groupe de tours est $\langle\ell\rangle\simeq
\mathbb Z$. Sur ce groupe, on déclare le caractère de signe
\begin{equation}
 \varepsilon:\langle\ell\rangle\longrightarrow\{+1,-1\},
 \qquad \varepsilon(\ell)=-1.
 \label{vi:vi:eq:holonomia-signo}
\end{equation}
Le groupe de tours enregistre combien de fois la boucle est parcourue ; son
type est différent de la phase finie $C_9$. Le fibré d’intervalles
associé au caractère est
\begin{equation}
 \mathcal M_\ell=([0,1]\times[-1,1])/
                  \bigl((0,u)\sim(1,-u)\bigr).
 \label{vi:vi:eq:banda-mobius}
\end{equation}

\begin{theorem}[Retour orienté sur la bande]
\label{vi:vi:thm:retorno-mobius}
Le fibré \eqref{vi:vi:eq:banda-mobius} est un ruban de Möbius.
Son transport sur $\ell$ change le signe de la fibre et sur
$\ell^2$ le restitue. Par conséquent, un vecteur non nul de cette droite
requiert deux tours pour son retour orienté. Cela conserve la
distinction entre retour d’orientation et retour de l’état complet.
\end{theorem}
\begin{proof}
La relation d’extrémités de \eqref{vi:vi:eq:banda-mobius} est l’application
linéaire $u\mapsto-u$, de déterminant négatif. Elle donne la droite réelle
non orientable sur le cercle et sa bande
d’intervalles correspondante. Composer deux transports produit $u\mapsto u$.
La mémoire de l’extension ne fait pas partie de ce quotient de signe ;
sa mise à jour continue d’obéir au transport du TPK et à
\eqref{vi:vi:eq:persistencia-fase-memoria}.
\end{proof}

Un paramétrage d’un tour de la base au moyen de $2\pi$
représente la restitution du revêtement au moyen de $4\pi$. Si le
lecteur temporel attribue $108$ pas au tour pertinent, le
relèvement orienté se ferme à $216$ pas. Ces attributions
requièrent la boucle et le lecteur indiqués ; elles ne divisent pas l’histoire en
deux particules successives de $108$ pas.

\begin{proposition}[Décomposition des observables sur le revêtement double]
\label{vi:vi:prop:observables-paridad}
Soit $p:\widetilde X\to X$ un revêtement principal double, d’involution
$\iota$. Pour $f\in C^0(\widetilde X,\mathbb R)$,
\[
 f^+=\frac{f+f\circ\iota}{2},\qquad
 f^-=\frac{f-f\circ\iota}{2}
\]
donnent $f=f^++f^-$, avec $f^+\circ\iota=f^+$ et
$f^-\circ\iota=-f^-$. La première partie descend à $X$ ; la
seconde représente une section de la droite de signe associée. Sur
la droite de Möbius, toute section globale continue possède un zéro.
\end{proposition}
\begin{proof}
Les identités s’obtiennent en substituant $\iota^2=I$. La partie
invariante est constante dans chaque fibre et définit une fonction sur le
quotient. La partie anti-invariante satisfait la loi de transition
de la représentation de signe. Sur le cercle, une section de cette
droite est représentée par une fonction réelle continue $u$ sur $[0,1]$
telle que $u(1)=-u(0)$. Si l’extrémité vaut zéro, il y a déjà un zéro ; dans le cas
contraire, les extrémités ont des signes opposés et le théorème des
valeurs intermédiaires fournit un zéro intérieur.
\end{proof}

La décomposition distingue deux classes d’observables sur un
même revêtement. Elle permet de représenter une information impaire d’orientation
et des données paires conservées. Pour l’appliquer à un registre matériel, on
utilise le relèvement \eqref{vi:vi:eq:levantamiento-cm} ; le caractère
de signe décrit la droite transportée, non toutes les données de la
particule.

\subsection{Réseau bidirectionnel et conjugaison des canaux}
\label{vi:vi:subsec:reticulo-bidireccional}

Le registre angulaire entier publie deux coordonnées
\begin{equation}
 W_+=6n_A+s_C,\qquad W_-=6n_A-s_C,
 \qquad (n_A,s_C)\in\mathbb Z^2.
 \label{vi:vi:eq:reticulo-ida-retorno}
\end{equation}
Les coefficients $n_A,s_C$ sont des lectures internes de route. Le facteur
$6$ appartient à la normalisation dodécaphasique de ce registre.

\begin{proposition}[Indice du réseau et échange des lectures]
\label{vi:vi:prop:indice-bidireccional}
L’image de \eqref{vi:vi:eq:reticulo-ida-retorno} a pour indice $12$
dans $\mathbb Z^2$ et pour forme normale de Smith $\operatorname{diag}(1,12)$.
Elle est caractérisée par $W_++W_-\equiv0\pmod{12}$. L’opération
$s_C\mapsto-s_C$ échange $W_+$ et $W_-$.
\end{proposition}
\begin{proof}
La matrice de l'application est
$\bigl(\begin{smallmatrix}6&1\\6&-1\end{smallmatrix}\bigr)$.
Le plus grand commun diviseur de ses entrées est $1$ et son déterminant
est $-12$, ce qui donne les facteurs invariants indiqués. Les
coordonnées inverses sont
$n_A=(W_++W_-)/12$ et $s_C=(W_+-W_-)/2$.
Si la somme est divisible par $12$, les deux entiers ont la même
parité et les deux expressions sont entières. La condition est également
nécessaire par la somme de \eqref{vi:vi:eq:reticulo-ida-retorno}.
L’affirmation sur l’échange se vérifie en substituant $-s_C$.
\end{proof}

Dans la représentation angulaire réelle, les coordonnées déjà engendrées
$a=A_{\mathrm{rad}}$ et $b=C^*_{\mathrm{rad}}$ produisent les
canaux $q_\pm=\exp[-(a\pm b)]$. Dans le domaine $a>|b|$, tous deux
appartiennent à $(0,1)$. Soit $R^*=R$, $R^2=I$ la graduation des
feuillets et soit $C$ un échange unitaire tel que $CRC^{-1}=-R$.
Pour $T(b)=\exp[-(aI+bR)]$, on obtient
\begin{equation}
 CT(b)C^{-1}=T(-b),\qquad
 CP_\pm C^{-1}=P_\mp,\quad P_\pm=(I\pm R)/2.
 \label{vi:vi:eq:conjugacion-canales}
\end{equation}
En effet, la conjugaison transforme l’exposant et chaque terme
de sa série convergente. Cette preuve transmet l’échange aux
canaux. Le réseau entier et l’opérateur réel ont des domaines
distincts et publient la même opération d’échange au moyen de
leurs applications de lecture déclarées.

\subsection{Représentation spinorielle du bloc électronique}
\label{vi:vi:subsec:spin-electron}

La construction APP tridimensionnelle des projecteurs de somme et de
produit sélectionne un plan principal électronique par le mode
tritique $(1,-1,0)^3$. Sa dérivation et l’évaluation des
projecteurs sont incluses parmi les antécédents du présent
article. Dans une base orthonormale de ce plan, leurs restrictions
sont
\begin{equation}
 P=\begin{pmatrix}1&0\\0&0\end{pmatrix},\qquad
 Q=\begin{pmatrix}1/4&\sqrt3/4\\\sqrt3/4&3/4\end{pmatrix},
 \qquad S=2P-I,\quad J=\frac4{\sqrt3}[P,Q].
 \label{vi:vi:eq:bloque-electronico}
\end{equation}
Cette représentation reçoit les deux projecteurs construits et
conserve leur orientation. La multiplication donne
$S^*=S$, $J^*=-J$, $S^2=I$, $J^2=-I$ et $SJ=-JS$.
Concrètement, $S=\operatorname{diag}(1,-1)$ et
$J=\bigl(\begin{smallmatrix}0&1\\-1&0\end{smallmatrix}\bigr)$.

Soit $\mathscr S_e=E_e\otimes_{\mathbb R}\mathbb C$ la
complexification du plan réel. L’unité scalaire $i$ appartient
au corps de cette représentation ; $J$ reste
l’endomorphisme réel provenant des projecteurs.

\begin{proposition}[Triplet hermitien et représentation des rotations]
\label{vi:vi:prop:spin}
Dans $\mathscr S_e$, définissons
\begin{equation}
 \sigma_1=SJ,\qquad \sigma_2=-iJ,\qquad \sigma_3=S.
 \label{vi:vi:eq:terna}
\end{equation}
Ces opérateurs sont hermitiens, de trace nulle, et satisfont
\begin{equation}
 \sigma_a\sigma_b=\delta_{ab}I+
 i\sum_c\epsilon_{abc}\sigma_c,
 \qquad L_a=\sigma_a/2,\qquad
 [L_a,L_b]=i\sum_c\epsilon_{abc}L_c,
 \quad\sum_aL_a^2=\tfrac34I.
 \label{vi:vi:eq:spin-algebra}
\end{equation}
Ils constituent la représentation irréductible de spin $1/2$ sur
le bloc électronique.
\end{proposition}
\begin{proof}
Les relations de $S,J$ donnent $(SJ)^*=SJ$, $(-iJ)^*=-iJ$ et
$\sigma_a^2=I$. Les produits cycliques sont
$(SJ)(-iJ)=iS$, $(-iJ)S=iSJ$ et $S(SJ)=J=i(-iJ)$ ;
inverser l’ordre change le signe. Cela prouve l’identité de
produit et les commutateurs. Dans la base de
\eqref{vi:vi:eq:bloque-electronico}, les trois matrices sont
\[
 \begin{pmatrix}0&1\\1&0\end{pmatrix},\quad
 \begin{pmatrix}0&-i\\i&0\end{pmatrix},\quad
 \begin{pmatrix}1&0\\0&-1\end{pmatrix}.
\]
Leur trace est nulle et, avec $I$, elles engendrent $M_2(\mathbb C)$.
Un sous-espace invariant par les trois est invariant par toute cette
algèbre et ne peut être que $0$ ou $\mathscr S_e$. Enfin,
sommer $L_a^2=I/4$ donne la valeur $3I/4$ de l’opérateur quadratique.
\end{proof}

La reconnaissance de ce triplet comme matrices de Pauli est postérieure
à \eqref{vi:vi:eq:bloque-electronico}. La norme et le produit
$\langle X,Y\rangle=\tfrac12\operatorname{tr}(XY)$ dans
l’espace réel des matrices hermitiennes de trace nulle donnent son espace
de directions. Pour $n\in\mathbb R^3$, $|n|=1$, soient
\begin{equation}
 H_e(n)=\sum_an_a\sigma_a,\qquad
 h_e(n)=\tfrac12H_e(n),\qquad
 \Pi_\pm(n)=\tfrac12(I\pm H_e(n)).
 \label{vi:vi:eq:helicidad}
\end{equation}

\begin{proposition}[Décomposition directionnelle et double retour]
\label{vi:vi:prop:helicidad}
Les opérateurs $\Pi_\pm(n)$ sont des projecteurs orthogonaux de
rang un et $h_e(n)\Pi_\pm(n)=\pm\tfrac12\Pi_\pm(n)$.
Le groupe de rotations autour de $n$ est
\begin{equation}
 U_n(\theta)=e^{-i\theta H_e(n)/2}
 =\cos(\theta/2)I-i\sin(\theta/2)H_e(n),\qquad
 U_n(2\pi)=-I,\quad U_n(4\pi)=I.
 \label{vi:vi:eq:retorno-spin}
\end{equation}
\end{proposition}
\begin{proof}
La partie antisymétrique de \eqref{vi:vi:eq:spin-algebra} s’annule
lorsqu’elle est contractée avec $n_an_b$, de sorte que $H_e(n)^2=I$.
Il en découle l’idempotence, l’orthogonalité et la somme
$\Pi_++\Pi_-=I$. Ils sont hermitiens et de trace un, donc leur
rang vaut un. Multiplier par $H_e(n)/2$ donne les valeurs propres.
Séparer les puissances paires et impaires de l’exponentielle produit
\eqref{vi:vi:eq:retorno-spin} et sa loi de groupe. La normalisation
angulaire se vérifie par
\[
 U_n(\theta)H_e(v)U_n(\theta)^*
 =H_e\bigl(v\cos\theta+(n\times v)\sin\theta
       +n(n\cdot v)(1-\cos\theta)\bigr),
\]
qui résulte de la multiplication avec l’identité de produit du
triplet. L’action induite sur les directions est une rotation
d’angle $\theta$ ; le relèvement spinoriel porte sa moitié.
\end{proof}

Si $n$ se réalise comme direction de propagation, $h_e(n)$ est
l’hélicité adimensionnelle. L’inversion de cette direction satisfait
$H_e(-n)=-H_e(n)$ et $\Pi_\pm(-n)=\Pi_\mp(n)$.
Cette opération change les étiquettes directionnelles ; la conjugaison
de charge transporte le registre matériel au moyen de
\eqref{vi:vi:eq:levantamiento-cm}. La chiralité se définit au moyen de
la graduation relativiste correspondant à la représentation et à
ses projecteurs, distincts de \eqref{vi:vi:eq:helicidad}.

L’évaluation scalaire électronique agit comme
$\varepsilon_e^\beta I$ sur $\mathscr S_e$ et commute avec
le triplet et avec $\Pi_\pm(n)$. L’information rotationnelle conserve
la structure de cette évaluation sans introduire d’autre facteur dans
la masse. De son côté, le ruban de Möbius transporte une droite
réelle autour d’une boucle et l’holonomie nonadique transporte
phase et mémoire. Les trois retours peuvent intervenir dans une
même réalisation ; leurs domaines et leurs opérateurs d’entrelacement déterminent la
composition précise.
